\documentclass[11pt,letterpaper]{article}
\usepackage[T1]{fontenc}
\usepackage{lmodern,amsmath,amssymb,booktabs}
\usepackage[left=21mm,right=21mm,top=18mm,bottom=18mm]{geometry}
\usepackage[hidelinks]{hyperref}
\hypersetup{pdftitle={Loop counts in closed meandric systems at large weight},pdfauthor={Iaroslav Argunov}}
\newcommand{\Prb}{\mathbb P}
\newcommand{\EE}{\mathbb E}
\newcommand{\Var}{\operatorname{Var}}
\newcommand{\Cov}{\operatorname{Cov}}
\newcommand{\Poi}{\operatorname{Pois}}
\newcommand{\TV}{d_{\mathrm{TV}}}
\newtheorem{theorem}{Theorem}
\newtheorem{lemma}{Lemma}
\raggedbottom
\linespread{0.96}
\setlength{\emergencystretch}{1em}
\renewcommand{\labelitemi}{\(\bullet\)}
\begin{document}
\begin{center}
{\LARGE Loop counts in closed meandric systems at large weight}\\[7pt]
{\large Finite boundary windows and a Poisson crossover}\\[8pt]
Iaroslav Argunov\\[3pt]
{\small 19 November 2026}
\end{center}
\begin{abstract}
Give each loop of a closed meandric system a factor \(q\). A finite-window
frontier estimate proves that the partition function has an
\(n^{-3/2}\) prefactor for every fixed real \(q\ge18\). The bound also
places the mean and fluctuations of the loop count in the range of a
conditional local central limit theorem. If \(q\) tends to infinity,
the defect from the maximal loop count is asymptotically Poisson in
total variation with an explicitly defined mean. Replacing that mean
by its leading term \(2n/q\) gives a sharp crossover at \(q\asymp
n^{1/3}\). The predicted transition of the partition function at
\(q=2\) remains outside the range certified here.
\end{abstract}

\section*{1. The composition and its saddle}
Choose an ordered pair of noncrossing perfect matchings above and below
\(2n\) linearly labelled points, counting a common arch as a loop on
two vertices. If the pair has \(k\) loops, its defect is \(D_n=n-k\).
Write \(M_n(y)=\sum y^{n-k}\) over all such pairs and \(m_n(q)=q^n
M_n(1/q)\). Under probability proportional to \(q^k\), \(D_n\) has
probability generating function \(M_n(s/q)/M_n(1/q)\).

Join the two matchings in the lattice of noncrossing partitions of
the \(2n\) points. A pair is irreducible if the join consists of one
block. Denote its generating function (including a constant term) by
\(F(z,y)=1+z+\sum_{a\ge2,s} f_{a,s}z^ay^s\), where \(s\) is defect.
The join construction and the irreducible support bound appear in
\href{https://arxiv.org/abs/1506.01345}{Nica, \S\S3--4} and
\href{https://doi.org/10.4171/AIHPD/80}{Fukuda and Nechita, Prop.~4.3
and Cor.~5.2}. Each nontrivial irreducible has \(a/2\le s<a\).
The block containing the first point leaves \(2a\) independent gaps
(including the interval after its last point). The restrictions to
these gaps can be chosen independently. The loops are disjoint across
the join blocks, so the defects add. Consequently, for
\(U(t,y)=1+\sum_{n\ge1}M_n(y)t^n\),
\[
 U=F(tU^2,y),\qquad
 M_n(y)=\frac1{2n+1}[z^n]F(z,y)^{2n+1}.
 \tag{1}
\]
The second equality is Lagrange inversion.

\begin{theorem}
For every \(0<y\le1/18\) there is a unique \(z=z(y)>0\) in the
interior of convergence of \(F\) such that \(zF_z/F=1/2\).
Give \((A,B)=(a,s)\) probability
\([z^ay^s]F\cdot z(y)^ay^s/F(z(y),y)\), using the coefficient of
the formal series in this expression. Put \(V=\Var A\). For each
fixed real \(q\ge18\), with \(y=1/q\) and \(z=z(y)\),
\[
 m_n(q)\sim \frac{F(z,y)}{4\sqrt{\pi V}}\,
 \frac{\{qF(z,y)^2/z\}^n}{n^{3/2}} .
 \tag{2}
\]
\end{theorem}
The exponent \(3/2\) in the large-weight phase of closed systems
was proposed by \href{https://arxiv.org/abs/cond-mat/0003008}{Di
Francesco, Guitter and Jacobsen, (2.38)}. We use finite windows of the boundary to establish the sufficient interval in Theorem 1.

\section*{2. A finite window at the river}
Write \(V(t,y)\) for the generating function of diagrams in \(U\) with
no arch common to both matchings. A common arch separates its endpoints
from the other vertices in the noncrossing join: its interior can have
no arch going to the exterior on either side. Conversely, a two-point
join block is a common arch. Removing the singleton term \(z\) in the
block decomposition therefore gives
\[
             V=F(tV^2,y)-tV^2.                                      \tag{3}
\]
We construct an upper bound on the weight of all continuations of a
partially drawn diagram. This supplies a lower bound on the convergence
radius of \(V\).

Scan the river vertices from left to right. A vertical cut meets two
lists of arch ends, one above and one below the river. Write each list
in order of increasing distance from the river. In the scanned half-plane
the ends are joined in pairs by disjoint paths. A length counts the
visited vertices on one path. The side alternates on crossing a vertex,
so opposite-side ends have odd length and same-side ends have even
length. Record a length by class 1 for length one, 2 for positive even
length, and 3 for odd length at least three. The four possible moves,
specified by the opening or closing parenthesis in each matching, are:
\begin{itemize}
\item \(U\): open on both sides, creating a path of length one with a
new nearest end on each list.
\item \(R,L\): close on one side and open on the other. Move the nearest
end of the closing list to the head of the other list and increment
its path length.
\item \(D\): close on both sides, consuming the nearest end on each
list. Two different paths of lengths \(i,j\) merge to length \(i+j+1\).
A shared path closes with weight \(q\); closure of a class-1 path is
forbidden because it creates a common arch.
\end{itemize}
A move requiring a missing end is unavailable. The two parentheses
at each vertex prescribe the exposed ends uniquely, by noncrossing
on each side. Consequently histories from empty to empty are in
bijection with the labelled arch pairs entering \(V\).

Let \(\mathcal S\) consist of all finite noncrossing pairings of these
ends with these parity-compatible classes. This includes every physically
reachable frontier. The four rules act on \(\mathcal S\) without actual
lengths: increment takes classes \(1,2,3\) to \(2,3,2\); a merge always
has length at least three and its parity is known. For positive weights
\(w_1,w_2,w_3\) put \(H(s)=\prod_p w_{c(p)}\) over open paths of \(s\).
For a valid move \(e:s\longrightarrow es\) let \(r_e(s)\) be
\(H(es)/H(s)\) times the closing weight, if any.

\begin{lemma}
Retain the nearest \(d\) ends of each list, or the entire list if
shorter. Also retain their partners and the classes of their paths.
Erase all other paths and compress the two lists, preserving order.
Call the resulting pairing \(P_ds\). For a move \(e\), its availability,
its multiplier, and \(P_d(es)\) are determined by \(P_{d+1}s\).
\end{lemma}
Indeed, at most one existing end is removed from either list in a single
move. Any old end that becomes one of its nearest \(d\) was among the
previous nearest \(d+1\) on its side (including an end about to slide
to the other side). A new end is given explicitly. The changed paths
meet the old nearest ends. All their old partners are therefore present,
including the far ends of two paths that merge. The order among
unaffected retained ends is preserved. If the full list had fewer than
\(d+1\) ends it was retained in its entirety. These facts also prove
\(P_dP_{d+1}=P_d\). This argument holds at arbitrary frontier size.

To enumerate the possible \(P_ds\), take all noncrossing pairings
of \(2j\) ends with \(0\le j\le 2d\), split the ordered ends at
every possible divider, and discard the pairing if some path misses
both near-end windows. The linear boundary order reverses the first
list and then follows the second. Assign class 2 to each same-side
path and either 1 or 3 to each opposite-side path. Every possible
window occurs in this list. Conversely each enumerated configuration
is its own window in \(\mathcal S\). For functions symmetric in the
two lists we identify a configuration with its reflection. Both slides
are still included in the sum, even if they give the same target.

Define the positive linear operator \(Tg(s)=\sum_e r_e(s)g(es)\).
A finite test for a multiplier bound \(\Lambda\) consists of positive
numbers \(h(P_ds)\) such that for each enumerated width-\((d+1)\)
window \(S\),
\[
       \sum_e r_e(S)\,h(P_d(eS))\ \le\ \Lambda h(P_dS).                \tag{4}
\]
The same \(S\) determines all terms on the left. The lemma gives (4)
on the whole state space. Iteration from the empty state bounds the
total closing weight of histories returning empty after \(2n\) steps by
\(\Lambda^{2n}\), since their initial and final potentials coincide.
Accordingly \(V(t,1/q)\) has radius at least \(q/\Lambda^2\).

\paragraph{What increasing the width can prove.}
For fixed \(q,w_i>0\) put \(a_j=\sup_{\mathcal S} T^j1\) and
\(\rho=\lim_{j\to\infty}a_j^{1/j}\). This limit exists by
submultiplicativity; each \(a_j\) is finite since the local multipliers
are bounded. The infima of the constants in (4) decrease with \(d\)
and converge to \(\rho\). In one direction any positive finite-window
potential has positive minimum and finite maximum. Iteration of (4)
thus gives \(a_j\le (\max h/\min h)\Lambda^j\). For the other direction
take \(\lambda>\rho\) and choose \(L\) such that \(a_L<\lambda^L\).
Each \(T^j1\) depends only on \(P_j s\), by the lemma. The positive
reflection-invariant function
\[
 g=\sum_{j=0}^{L-1}\lambda^{-j}T^j1
 \qquad\hbox{satisfies}\qquad
 Tg=\lambda(g-1+\lambda^{-L}T^L1)\le\lambda g .
\]
It is a width-\((L-1)\) certificate. The monotonicity uses
\(P_dP_{d+1}=P_d\). Here \(\rho\) is the uniform growth of continuations
from arbitrary formal states. Identifying it with the empty-to-empty
growth rate is a separate question.

\paragraph{The rational certificate.}
Take \(q_0=18\), \(d=6\), \((w_1,w_2,w_3)=(1432,2940,4284)/1000\),
and \(\Lambda=5447/1000\). Up to reflection there are 25,000 width-six
windows and 128,600 width-seven extensions. The certificate enumerates
them by matching the first end of each interval to an admissible partner
and splitting into its interior and exterior subintervals. A pair with
both endpoints outside the retained window is excluded at its creation.
A separate count by reflection and Burnside's lemma verifies both totals.
Start every width-six entry at \(10^{10}\). For each source take the
maximum over its extensions of the left side of (4). Normalize the new
vector by its maximum, multiply by \(10^{10}\), round down and replace
entries below one by one. After 180 iterations the largest exact ratio
over the extensions is
\[
   299370707754862/55041444031875 < 5.447 .
\]
Only the final inequalities are needed for the bound.

Set \(t_0=6065/10000\) and \(b=6977/10000\); then
\(t_0\Lambda^2<18\). For a lower bound let \(K_a\) be the sum of
\(18^{\#\text{loops}}\) over empty-to-empty histories of length \(2a\)
with at most 14 active ends at every stage and no common arch.
The forward recurrence starts at the empty state with weight one.
It drops moves exceeding the cap, assigns multiplier 18 to an allowed
closure, and sums incoming weights at the next state. All other
multipliers are one. It visits 18,924 reflected states and gives
\[
        t_0\left(\sum_{a=0}^{40}K_a(t_0/18)^a\right)^2 > b .
\]
Thus \(z_0=t_0V(t_0,1/18)^2>b\), and (3) implies convergence of
\(F(z,1/18)\) beyond \(b\). As a second check, differentiating (3)
at \(z=tV^2\) gives
\[
 (2zF_z-F)(V+2tV_t)=V^2(tV-1+2t^2V_t).
\]
The same restricted counts give a lower bound greater than 0.00157 for
\(t_0V-1+2t_0^2V_t\). The derivatives exist at this point because
\(t_0\) is strictly inside the radius and \(tV(t,1/18)^2\) is increasing.

For \(q=18x^2\), \(x\ge1\), the support inequality \(2s\ge a\) for all
nontrivial blocks now implies convergence of \(F(z,1/q)\) beyond \(bx\).
At that argument a lower bound for \(2zF_z-F\) is
\[
 bx-1+\sum_{2\le a\le12,s} C_{a,s}x^{a-2s},\qquad
 C_{a,s}=(2a-1)f_{a,s} b^a18^{-s}.
\]
The coefficients through order twelve follow by reverting (1) from the
complete loop-count rows. Exact arithmetic gives a value greater than
0.00016 at \(x=1\) and derivative at least
\(b-\sum(2s-a)C_{a,s}>0.4948\) throughout \(x\ge1\).
The mean \(zF_z/F\) starts at zero and has strictly positive derivative
\(\Var A\) with respect to \(\log z\). This establishes the interior
saddle throughout the claimed interval. At positive \(y\), convergence
extends to a neighbourhood in both variables since \(s\le a\).
Near \((z,y)=(1,0)\), the bounds \(f_{a,s}\le16^a\) and \(s\ge a/2\)
give an analytic neighbourhood. The saddle there is analytic by the
implicit function theorem applied to \(2zF_z-F\).


\section*{3. The conditional law}
Take \(m=2n+1\) independent copies of \((A,B)\) at the saddle.
Let \(S_A,S_B\) be their two sums. Formula (1) yields exactly
\[
 \Prb(D_n=r)=\Prb(S_B=r\mid S_A=n),\qquad
 M_n(y)=F(z,y)^m z^{-n}\Prb(S_A=n)/m.
 \tag{5}
\]
Write \(\beta=\EE B,\ c=\Cov(A,B),\ W_*=\Var B-c^2/V\), and set
\(\eta=m\beta-c/(2V)\). At fixed positive \(y\), the atoms
\((0,0),(1,0),(2,1)\) give \(W_*>0\) and lattice span one.
Implicit differentiation of the saddle equation gives
\[
 \frac{d\beta}{d\log y}=W_*,\qquad
 2\beta=2y+4y^2-16y^3+O(y^4),\qquad W_*=y+O(y^2).
 \tag{6}
\]
Indeed \(d\log z/d\log y=-c/V\). Through degree three the marked
terms of \(F\) are \(2z^2y+(8z^3+4z^4)y^2+
(42z^4+60z^5+18z^6)y^3\). Substitution into \(2zF_z=F\)
gives \(z=1-6y+4y^2+96y^3+O(y^4)\); then evaluate \(2yF_y/F\).
\begin{theorem}
Let \(q=q_n\ge18\) and \(ny=n/q\to\infty\). Uniformly over
integers \(r\),
\[
 \sqrt{2\pi mW_*}\Prb(D_n=r)
 =\exp[-(r-\eta)^2/(2mW_*)]+o(1).
 \tag{7}
\]
In particular \((D_n-\eta)/\sqrt{mW_*}\) tends to a standard normal.
For fixed \(q\ge18\) the proportion of loops tends in probability
to \(1-2\beta\).
\end{theorem}
General conditional Gaussian theorems for triangular arrays are given
by \href{https://doi.org/10.1239/jap/996986753}{Janson}.
We include the local argument because \(y\) can decrease to zero.
Over \(0<y\le1/18\) the saddle varies in a compact positive range,
with a uniform analytic neighbourhood at each limit point after
including \(y=0\). The first two atoms above thus have probabilities
bounded below, the third probability at least a constant times \(y\).
The third centered absolute moments of \(A,B\) are \(O(1)\) and
\(O(y)\). For their characteristic function \(\varphi\),
\[
 |\varphi(u,v)|^m\le
 \exp\{-cm(\|u\|^2+y\|2u+v\|^2)\},
\]
where here \(c>0\) is a constant and the norms are distances to
\(2\pi\mathbb Z\). This follows by retaining the two independent
differences from \((0,0)\) in \(1-|\varphi|^2\).
The torus change \((u,v)\mapsto(u,2u+v)\) has determinant one.
Outside the rectangle with scaled coordinates bounded by \(L\),
the integral is \(O(e^{-c'L^2}/(m\sqrt y))\), with respective
scales \(1/\sqrt m\) and \(1/\sqrt{my}\).
Inside, the error in the centered Gaussian log expansion is
\(O(m|u|^3+my|v|^3)=o(1)\) at fixed \(L\). The covariance form
is bounded above and below by constant multiples of \(u^2+yv^2\).
These estimates give absolute error \(o(1/(m\sqrt y))\) in the
bivariate normal lattice approximation, uniformly in the target
because its Fourier phase has modulus one. The one-dimensional
argument for \(S_A=n\) gives \((2\pi mV)^{-1/2}(1+o(1))\);
the discrepancy \(n-m/2=-1/2\) is bounded. This last estimate
is uniform even without the requirement \(my\to\infty\), since
the two atoms of \(A\) at 0 and 1 remain bounded below.
Equation (5) now gives (2), uniformly in this latter relative
estimate for the corresponding parameters. Dividing the normal
densities gives exactly the conditional mean \(\eta\) and variance
\(mW_*\) in (7). Summing on fixed standardized intervals and
then increasing their lengths proves the weak limit.

\section*{4. Which Poisson mean?}
Here \(\TV\) is half the sum of absolute differences of point
probabilities, and \(\Phi\) denotes the standard normal distribution
function.
\begin{theorem}
Along every sequence \(q=q_n\to\infty\),
\[
 \TV\big(\mathcal L(D_n),\Poi(m\beta)\big)\longrightarrow0.
\]
If \(n/q\to\lambda<\infty\), the limit law is \(\Poi(2\lambda)\),
including a point mass at zero when \(\lambda=0\). With the
simpler mean \(2n/q\), the limiting distance is
\[
 \lim\TV\big(\mathcal L(D_n),\Poi(2n/q)\big)=
 \begin{cases}
 0, & q/n^{1/3}\to\infty,\\
 2\Phi(\sqrt2/a^{3/2})-1,& q/n^{1/3}\to a\in(0,\infty),\\
 1,& q/n^{1/3}\to0.
 \end{cases}
 \tag{8}
\]
\end{theorem}
For bounded \(ny\to\lambda\), (5) and the uniform one-dimensional
estimate give \(M_n(y)/C_n\to e^{2\lambda}\), with \(C_n\) the
\(n\)-th Catalan number. Indeed \(V\to1/4,\ F\to2\),
\(F^2/(4z)=1+2y+O(y^2)\). For fixed \(r\), (1) gives
\([y^r]M_n(y)/C_n\sim(2n)^r/r!\). One way to see this is to expand
\(F(z,y)^m\): the term using \(r\) distinct copies of \(2z^2y\)
is \(2^r\binom mr\binom{m-r}{n-2r}/m\) before normalization.
Every other term selects at most \(r-1\) nonconstant-in-\(y\)
factors of bounded degree. Dividing by \(C_n=\binom{m}{n}/m\)
gives the stated equivalent. Pointwise convergence of masses to a
probability on the integers implies convergence in total variation.

When \(ny\to\infty\), (7) and the Poisson local limit theorem
(equivalently Stirling's formula on fixed central windows) compare
the two lattice laws on intervals of length \(O(\sqrt{ny})\).
Their masses outside increasing such windows vanish.
For mean \(m\beta\) the normal variances have ratio tending to one,
and the two centers differ by \(O(1)\). This proves the first claim
also in this range and hence by subsequences for all \(q\to\infty\).
For mean \(2ny\), (6) instead displaces the centers by
\[
 \frac{\eta-2ny}{\sqrt{2ny}}
 = 2\sqrt{2}\sqrt{ny^3}(1+O(y))+o(1).
\]
When this has a finite limit \(d\), summing the absolute differences
by Riemann sums gives the distance between two unit normals separated
by \(d\), equal to \(2\Phi(d/2)-1\). When the displacement diverges,
the central windows eventually separate and the distance tends to one.
This proves (8). The cube-root scale concerns the use of a particular
truncated mean; the large-\(q\) exponent in (2) addresses a different
limit.

The first correction to the Poisson mean gives a second explicit
scale. For \(q\to\infty\) and mean \(2n/q+4n/q^2\), the
distance tends to zero for \(q/n^{1/5}\to\infty\), and to one for
\(q/n^{1/5}\to0\). At \(q/n^{1/5}\to a\in(0,\infty)\)
its limit is \(2\Phi(4\sqrt2/a^{5/2})-1\). In fact (6) makes the
new standardized displacement \(-8\sqrt2\sqrt{ny^5}(1+O(y))+o(1)\)
when \(ny\to\infty\), so the same comparison proves all three limits.
If \(ny\) is bounded the two means differ by \(o(1)\).

\paragraph{Reproducibility.}
The companion certificate uses the Python standard library and exact arithmetic. It tests the transition
majorant against explicit port pairings and path lengths and
recomputes all finite series entries from pair histories.
The location of a sharp transition in the fixed-weight model and
the exponent for connected meanders remain open.

\section*{References}\small
\noindent
P. Di Francesco, O. Golinelli and E. Guitter,
\href{https://arxiv.org/abs/hep-th/9607039}{Meanders: A direct
enumeration approach}, Nuclear Physics B 482 (1996), 497--535.
Their large-weight expansions provide the enumerative background. The discussion after (5.6) in their preprint distinguishes the limit at infinite weight from the exponent at fixed finite weight.\\[4pt]
P. Di Francesco, E. Guitter and J. L. Jacobsen,
\href{https://arxiv.org/abs/cond-mat/0003008}{Exact meander
asymptotics: a numerical check}, Nuclear Physics B 580 (2000), 757--795.\\[4pt]
M. Fukuda and I. Nechita,
\href{https://arxiv.org/abs/1609.02756}{Enumerating meandric systems
with large number of loops}, Ann. Inst. Henri Poincar\'e D 6
(2019), 607--640.\\[4pt]
S. Janson, \href{https://doi.org/10.1239/jap/996986753}{Moment
convergence in conditional limit theorems}, Journal of Applied Probability
38 (2001), 421--437.\\[4pt]
I. Jensen, \href{https://arxiv.org/abs/cond-mat/0008178}{A transfer
matrix approach to the enumeration of plane meanders}, Journal of
Physics A 33 (2000), 5953--5963.\\[4pt]
A. Nica, \href{https://arxiv.org/abs/1506.01345}{Free probability
aspect of irreducible meander systems, and some related observations
about meanders}, Infinite Dimensional Analysis, Quantum Probability
and Related Topics 19 (2016), \href{https://doi.org/10.1142/S0219025716500119}{DOI 10.1142/S0219025716500119}.
\end{document}
